% Groente- en fruitkalender
\cleardoublepage
\par\noindent{\handfont\fontsize{60}{58}\selectfont\color{terracotta}Seizoens-\\kalender\par}
 \vspace{6pt}{\color{terracotta}\rule{2.6cm}{0.5pt}\;\raisebox{-0.2ex}{$\scriptstyle\blacklozenge$}}\par
 {\vspace{6pt}\itshape\small\color{inkfaint}Wanneer onze groente en al het fruit op hun best zijn.\par}\vspace{1.2em}

Hieronder staat per groente en fruit het Nederlandse teeltseizoen. Bij de
groente hebben we ons beperkt tot wat in dit boek ook echt gebruikt wordt, zo
blijft de lijst een praktische geheugensteun in plaats van een volledige
opsomming. Fruit eten we het hele jaar door los, dus daarvan staat de
complete lijst erbij.

% Local, appendix-only helper: keeps the 45 near-identical entries below from
% being hand-formatted one by one.
\newcommand{\seizoenitem}[2]{\noindent#1: #2\par\vspace{0.4em}}

\blockrule{Groente}
\begin{multicols}{2}
\small
\seizoenitem{Andijvie}{april tot en met december}
\seizoenitem{Artisjok}{juni tot en met oktober}
\seizoenitem{Asperge}{april tot en met juni}
\seizoenitem{Aubergine}{juli tot en met oktober}
\seizoenitem{Bleekselderij}{juli tot en met november}
\seizoenitem{Bloemkool}{maart tot en met november}
\seizoenitem{Boerenkool}{januari tot en met maart, oktober tot en met december}
\seizoenitem{Broccoli}{juni tot en met november}
\seizoenitem{Courgette}{juli tot en met november}
\seizoenitem{Doperwt}{mei tot en met juli}
\seizoenitem{Komkommer}{juli tot en met oktober}
\seizoenitem{Maïs}{augustus tot en met september}
\seizoenitem{Paddenstoelen}{het hele jaar door}
\seizoenitem{Paprika}{juli tot en met oktober}
\seizoenitem{Pastinaak}{januari tot en met maart, oktober tot en met december}
\seizoenitem{Peultjes}{mei tot en met juli}
\seizoenitem{Pompoen}{januari tot en met april, augustus tot en met december}
\seizoenitem{Prei}{het hele jaar door}
\seizoenitem{Sperzieboon}{juni tot en met oktober}
\seizoenitem{Spinazie}{maart tot en met oktober}
\seizoenitem{Spitskool}{mei tot en met oktober}
\seizoenitem{Spruiten}{januari tot en met maart, oktober tot en met december}
\seizoenitem{Tomaat}{juni tot en met oktober}
\seizoenitem{Ui}{het hele jaar door}
\seizoenitem{Venkel}{juni tot en met november}
\seizoenitem{Witlof}{januari tot en met april, oktober tot en met december}
\seizoenitem{Wortel}{het hele jaar door}
\end{multicols}

\blockrule{Fruit}
\begin{multicols}{2}
\small
\seizoenitem{Aardbei}{april tot en met augustus}
\seizoenitem{Abrikoos}{juni tot en met augustus}
\seizoenitem{Appel}{januari tot en met mei, augustus tot en met december}
\seizoenitem{Blauwe bes}{juni tot en met september}
\seizoenitem{Braam}{juli tot en met oktober}
\seizoenitem{Druif}{augustus tot en met november}
\seizoenitem{Framboos}{juni tot en met oktober}
\seizoenitem{Kers}{juni tot en met augustus}
\seizoenitem{Kiwibes}{september tot en met oktober}
\seizoenitem{Kweepeer}{september tot en met oktober}
\seizoenitem{Meloen}{augustus tot en met november}
\seizoenitem{Nectarine}{juni tot en met september}
\seizoenitem{Peer}{januari tot en met mei, augustus tot en met december}
\seizoenitem{Perzik}{juni tot en met september}
\seizoenitem{Pruim}{augustus tot en met september}
\seizoenitem{Rode bes}{juni tot en met juli}
\seizoenitem{Vijg}{augustus tot en met september}
\seizoenitem{Zwarte bes}{juni tot en met augustus}
\end{multicols}

\vspace{0.8em}
{\footnotesize\itshape\color{terracotta}Bron: Voedingscentrum, \emph{Seizoengroente- en
fruitkalender}. \url{https://www.voedingscentrum.nl/Assets/Uploads/voedingscentrum/Documents/Consumenten/Veelgestelde%20vragen/Voedingscentrum%20seizoengroente-%20en%20fruitkalender.pdf}\par}
